%%%%%%%%%%%%%%%%%%%%%%%%%%%%%%%%%%%%%%%%%%%%%%%%%%%%%
% Folienvorlage fuer Latex Beamer / PDF Praesentation
% Vorlage erstellt durch
% Alexander Kr�ger 
% 3.4.2008
%%%%%%%%%%%%%%%%%%%%%%%%%%%%%%%%%%%%%%%%%%%%%%%%%%%%%

\documentclass[8pt]{beamer}
%%%%%% Choose theme of frames. %%%%%%%%%%%%%%%%%%%%%%%%%
%%%%%% Note that there are many more of them than those given here. %%%%

% Navigation panel 
%\usetheme{Warsaw}
% Miniframe navigation panel
\usetheme{Darmstadt}
% no navigation panel
%\usetheme{Madrid}

%%%%%%%%%%%%%%%%%%%%%%%%%%%%%%%%%

%%% Choose size of frames
% \geometry{papersize={12.8cm,15.5cm}}
%%%%%%%%%%%%%%%%%%%%%%%%%%%%%%%%%%%%%%

% suppress navigation panel for navigation through frames
\setbeamertemplate{navigation symbols}{}
%%%%%%%%%%%%%%%%%%%%%%%%%%%%%%%%%%%%%%%%

\usepackage[english]{babel}
\usepackage{times}
\usepackage[T1]{fontenc}

\usepackage{amssymb}
\usepackage{subfigure}
\usepackage{graphicx}
\usepackage{psfrag}
\usepackage{pgf}

\usepackage{xcolor}

\itemindent 0cm

% Insert of logo on all frames
% placement is determined by chosen frame
\pgfdeclareimage[width = 1.15\textwidth]{logo}{Images/logo_fgnt_und_upb2}
\logo{\pgfuseimage{logo}} 
%%%%%%%%%%%%%%%%%%%%%%%%%%%%%%%%%%%%%%%%%%%%%%%%%%%%%%%%%%%%%%%%%%%%%%%%%%%%%%%%%%%%%%%%%%%%

% Set position of logo manually. %%%%%%%
% Logo has to be inserted on each frame by the command \MyLogo manually. %%%%%%%%%%%%%%%%%%
%\usepackage[absolute,overlay]{textpos}
%\setlength{\TPHorizModule}{1.0\textwidth}
%\setlength{\TPVertModule}{1.0\textheight}
%\newcommand{\MyLogo}{%
%\begin{textblock}{1}(0.0,0.92)
%  %\pgfuseimage{logo}
%  \includegraphics[width = 1.15\textwidth]{Images/logo_fgnt_und_upb.eps}
% \end{textblock}
%}
%%%%%%%%%%%%%%%%%%%%%%%%%%%%%%%%%%%%%%%%%%%%%%%%%%%

% Insert a logo on title page. %%%%%%%%%%%%%%%%%%%%%
% Position depends on chosen theme. %%%%%%%%%%%%%%%
%\pgfdeclareimage[width = 0.6\textwidth]{titellogo}{Images/logo_fgnt_fak_englisch}
%\titlegraphic{\pgfuseimage{titellogo}}
%%%%%%%%%%%%%%%%%%%%%%%%%%%%%%%%%%%%%%%%%%%%%%%%%%%%%%%%%%%%%%%%%%%%%%%%%%%%%%%%%%

% Mathematic command declarations %%%%%%%%%%%%%%%%%%%%%%%%%%%%%%%%%%%%%%%%%%%%%%%%%%%%%%%%%%%%%%%%%%%%%%%%%%%%%%%%%%%%%%%%%
	\newcommand{\grad}{\ensuremath{^\circ}}
	%define argmax function
	\DeclareMathOperator*{\argmax}{argmax}
	\newcommand{\vect}[1]{\ensuremath{\mathbf{#1}}}	
	\newcommand{\FFT}{\hspace{0.3cm}{\circ}\hspace{-.4mm}{\text{---}}{\hspace{-1mm}}{\bullet}\hspace{0.3cm}}
%%%%%%%%%%%%%%%%%%%%%%%%%%%%%%%%%%%%%%%%%%%%%%%%%%%%%%%%%%%%%%%%%%%%%%%%%%%%%%%%%%%%%%%%%%%%%%%%%%%%%%%%%%%%%%%%%%%%%%%%%%%%%

% Modify footline to show page numbers in Warsaw theme
\setbeamertemplate{footline}
{%
\begin{beamercolorbox}[wd=0.5\textwidth,ht=3ex,dp=1.5ex,leftskip=.5em,rightskip=.5em]{author in head/foot}%
\usebeamerfont{author in head/foot}%
\hfill\insertshortauthor\enspace%
\end{beamercolorbox}%
\vspace*{-4.5ex}\hspace*{0.5\textwidth}%
\begin{beamercolorbox}[wd=0.6\textwidth,ht=3ex,dp=1.5ex,left,leftskip=.5em]{title in head/foot}%
\usebeamerfont{title in head/foot}%
\enspace\insertshorttitle\hfill\insertframenumber{}/\inserttotalframenumber
\end{beamercolorbox}%
}
%%%%%%%%%%%%%%%%%%%%%%%%%%%%%%%%%%%%%%%%%%%%%%%%%%%%%%%%%%%%%%%%%%%%%%%%

\title[Short title] 
{Statische Sprachmodellierung }

\subtitle{Training Sprachmodellen in SLM-Toolkit}

\institute{
%\includegraphics[width=0.4 \textwidth]{Images/logo_fgnt_fak_englisch} \\ ~ \\ 
%University of Paderborn
}

\author[Author (short)] 
{Buyu Xiao}

\date{Date}

\subject{Subject}



\begin{document}

% title frame
\begin{frame}
  \titlepage
\end{frame}



% outline frame
\begin{frame}
  \frametitle{Outline}
  \tableofcontents
\end{frame}

\section{Section}

\subsection{Subsection}


% outline frame
\begin{frame}
  \frametitle{Outline}
  \tableofcontents
\end{frame}

\section{Sprach Modell}

\begin{frame}
  \frametitle{Sprach Modell}
  \tableofcontents
\end{frame}

%Statische Sprach_Modell

Was ist ein statisches Sprachmodell? Ein statische Sprachmodell wird durch die Wahrscheinlichkeit $p(w_{1}^n)$ von Wortfolgen  dargestellt.
Im folgenden Anwendungsbereich[1] spielt das Sprachmodell eine grosse Rolle.
\begin{itemize}
	\item Vereinfachen eines Spracherkenners
	\item Text-Komprimierung
	\item Extraktion von Schl\"usselw\"ortern aus Texten
	\item etc...
\end{itemize}
Im Detail, extrahiert das Sprachmodell die Eigenschaft, das Schlusswort oder die Syntax aus der Sprache.\footnote{Eine Sprache hat riesiggro\ss en Wortschatz und unendliche M\"oglichkeiten der Wortfolgen. D.h. ein Spracherkenner soll unendliche Wortfolgen kennen. Aber eine Sprache ist immer regul\"ar. Durch das Sprachmodell muss ein Spracherkenner  nur  limitierte Regeln speichern.}
Bei der Text-Komprimierung speichert man nur einmal die im Text immer wieder auftretenden W\"orter. Ein bekanntes Beispiel bei der Extraktion von Schl\"usselw\"orter ist das Tippprogramm. Man gibt erst Buchstaben ein und das Programm zeigt die m\"oglichen \"ubrigen Teile.

\section{Example beim SLM-Toolkit}
\subsection{Subsection}
\section{Zussammenfassung}




% frames of talk
\begin{frame}
  \frametitle{Example}

  \begin{block}{Block}
    This is a block.
  \end{block}

  \begin{exampleblock}{Example block}
    This is an example block.
  \end{exampleblock}

  \begin{alertblock}{Alert block}
    This is an alert block.
  \end{alertblock}
\end{frame}

\end{document}